\scanpage{70}
\mspage{23}
\deleted{2}4. \uline{Topos canoniques associés} à \uline{une stratification
globale.}

On va montrer comment, à une situation
stratifiée donnée, on peut en associer d'autres.

A) Images inverses générales. Rappelons les
axiomes utilisés jusqu'à présent :
\[
(1)\quad\left\{\begin{aligned}
&\text{a) }X_i\text{ fermés},\ (X_i)_{i\in I}\text{ loc. \unsure{finie|fini}}\\
&\text{b) }i\leq j\Longrightarrow X_i\leq X_j\\
&\text{c) }X=\bigcup X_i\\
&\text{d) }\forall i,j\in I,\ \unsure{\text{on a}|\ldots}\quad X_i\cap X_j=\bigcup_{\uline{\unsure{h|k}}<i,j}X_k
\end{aligned}\right.
\]

Notons que pour tt $X'$ au-dessus de
$X$, la famille des
\[
(2)\qquad X'_i\overset{\mathrm{def}}=\deleted{\unsure{X'\times_X}}X_i\times_X X'
\]
satisfait les \illegible{1} $=^{\mathrm s}$ conditions.

D'ailleurs le système $\frac12$ simpl des
$X_{\Delta_r}$ \unsure{a|—} comme image inverse celui
des $X'_{\Delta_r}$, défini par les $X'_i$,
\deleted{\unsure{d}} i.e. des iso
\[
(3)\quad\left\{\begin{aligned}
X'_{\Delta_r}&\simeq X_{\Delta_r}\times_X X'\\
X'^{*}_{\Delta_r}&\simeq X^*_{\Delta_r}\times_X X'\\
\dot X'_{\Delta_r}&=\dot X_{\Delta_r}\times_X X'
\end{aligned}\right.
\]

\margin{69}
